Wie du dem folgenden Stundenplan entnehmen kannst, enthält das Bioinformatik-Studium im ersten
Semester neben den Informatik- und Biochemie-Vorlesungen auch eine gute Portion Mathe.
%\vspace{-1cm}
\begin{center}
    \resizebox{\textwidth}{!}{% <- % is important
        \begin{tabular}{|c|c|c|c|c|c|}
        \hline
        Zeit     & Montag        & Dienstag                 & Mittwoch      & Donnerstag               & Freitag            \\ \hline
        08 -- 09 &               &                          &               &                          &                    \\ \hline
        09 -- 10 &               &                          &               &                          &                    \\ \hline
        10 -- 11 & Mathematik II &                          & Mathematik II &                          & Biochemie          \\ \cline{1-1} \cline{3-3} \cline{5-5}
        11 -- 12 & (\Matheprof)  &                          & (\Matheprof)  & Biochemie                & (\Biochemieprof)   \\ \hline
        12 -- 13 &               &                          &               &                          &                    \\ \cline{1-3} \cline{5-6}
        13 -- 14 &               &                          & Mathe         &Einf. in die Bioinformatik\footnotemark &                    \\ \cline{1-3} \cline{5-6}
        14 -- 15 &               & Praktische Informatik II & Rechenzentrum & Praktische Informatik II &                    \\ \cline{1-2} \cline{6-6}
        15 -- 16 &               & (\Infoprof)              & (ab 24.04.)   & (\Infoprof)              &                    \\ \cline{1-3} \cline{5-6}
        16 -- 17 &               &                          &               &                          &                    \\ \cline{1-3} \cline{5-6}
        17 -- 18 &               &                          &               &                          &                    \\ \hline
        %18 -- 19 &               &                          &               &                          &                    \\ \cline{1-1} \cline{3-6}
        %19 -- 20 &               &                          &               &                          &                    \\ \hline
        \end{tabular}% <- this one aswell
    }
\end{center}
\footnotetext{Einführung in die Bioinformatik (Prof. Thiel)}
\noindent Dieser Plan gilt für das zweite Semester Bioinformatik, für euch ist es jedoch das Erste.\\
